\documentclass{article}
\usepackage{graphicx} 
\usepackage{xeCJK}
\usepackage{setspace} 
\usepackage{type1cm}
\usepackage{indentfirst} 
\setlength\parindent{24pt} 
\singlespacing
\usepackage[top=1.27cm,bottom=1.27cm,left=1.27cm,right=1.27cm,a4paper]{geometry}
\setCJKmainfont{標楷體.ttc} %切換字體照左方名稱打

\begin{document}

\begin{center}
    \begin{LARGE}
    題目名稱
    \end{LARGE}
\end{center}

%%%%%%%%%%註解範例%%%%%%%%%%%%
\section{註解範例}
\begin{large}

    註解放在\footnote{註解}

    %範例
    馬克是個大好人，他常常被妹子\footnote{辭典檢視[妹子: ㄇㄟˋ ˙ㄗ] - 教育部《重編國語辭典修訂本》2021}告白
    
    %十二、參考資料來源之標註，必須於該引用文字結尾處，以註腳方式標示於該頁下方，需清楚標示參考資料之作者、文章名稱、期刊名稱或書籍名稱或網頁網址、期刊卷期或書籍章節、頁次、年份等資訊。
\end{large}
%%%%%%%%%%註解範例%%%%%%%%%%%%

%%%%%%%%%%圖表範例%%%%%%%%%%%%
\begin{figure}[h]
    \centering
    \centerline{\includegraphics[scale=0.1]{圖.png}} 
    \begin{center}
    \fontsize{12pt}
        0這是圖表
    \end{center}
\end{figure}
%%%%%%%%%%圖表範例%%%%%%%%%%%%


%%%%%%%%%%%%%%這是%%%%%%%%%%%%
%%%%%%%%%%%%%%內文%%%%%%%%%%%%

\begin{center}
\fontsize{18pt}
     0賽題研析
\end{center}
\begin{large}
    內文(2pages)
\end{large}

\begin{center}
\fontsize{18pt}
     0科學原理
\end{center}
\begin{large}
    內文(3pages)
\end{large}

%%%%%%%%%%%%%%%%%%%20pages%%%%%%%%%%%%%%%%%%%
\begin{center}
\fontsize{18pt}
     0實驗設計與規劃
\end{center}
\begin{large}
    內文
\end{large}


\begin{center}
\fontsize{18pt}
     0實驗結果與分析
\end{center}
\begin{large}
    內文
\end{large}


\begin{center}
\fontsize{18pt}
     0誤差來源與分析
\end{center}
\begin{large}
    內文
\end{large}


\begin{center}
\fontsize{18pt}
     0結論
\end{center}
\begin{large}
    內文
\end{large}

%%%%%%%%%%%%%%%%%%%20pages%%%%%%%%%%%%%%%%%%%

\begin{center}
\fontsize{18pt}
     0實驗心得
\end{center}
\begin{large}
    內文(1page)
\end{large}

%%%%%%%%%%%%%%這是%%%%%%%%%%%%
%%%%%%%%%%%%%%內文%%%%%%%%%%%%


\end{document}


%看不懂自己上網查不要來問我我會關通知
%https://www.cs.pu.edu.tw/~wckuo/doc/latex123/node5.html